%% appendix-mortal -- the constants and the proofs for the mortal case:
%% the separator word and the Pi-reader, which admits self events; positive
%% desubstitution; anchored exchange types, normal form and bounded
%% dictionaries; the unary top level; agreement with channels; positive column
%% summaries (S_ell has no F-term); supported compression; the proof of the
%% anchored compression theorem; the closed-form estimates and the decision
%% procedures.

\section{Proofs for Mortal Holes}\label{app:mo}

This appendix defines the constants of \cref{thm:mo-anchored} and proves the
results of \cref{sec:mo} whose proofs were omitted or outlined there.
\Cref{app:mo-separator} proves \cref{lem:mo-sep,lem:mo-pi-reader}, and
\cref{app:mo-desub} proves
\cref{lem:mo-desub}. \Cref{app:mo-exchange,app:mo-unary,app:mo-agreement,%
app:mo-columns,app:mo-compression,app:mo-proof} prove \cref{thm:mo-anchored},
following the structure of the proof of \cref{thm:bg-search} by \citet{Zhao26} as
recalled in \cref{app:bg-search}; each subsection states what it reuses and
what it adds. \Cref{app:mo-estimates} proves the closed-form estimates and
\cref{cor:mo-decidable}.

What is reused unchanged: the alignment lemma (\cref{lem:bg-alignment}); the
label facts \cref{eq:bg-label-facts} and the desubstitution lemma
(\cref{lem:bg-desub}); the facts on ordinary and symbolic states, start
prefixes and entry states (\cref{app:bg-monitors,app:bg-search}); the
comparison-preserving map lemma, the shape of level maps and the isolated
exchange for holes (\cref{lem:bg-comparison-map,lem:bg-level-maps,%
lem:bg-isolated-exchange}); the propagation of boundary pairs
(\cref{lem:bg-dictionary}(2)); the shortening of unary holes
(\cref{lem:bg-unary}); and the column summaries, column
shortening and agreement under derivation
(\cref{lem:bg-columns,lem:bg-agreement}). What is new is that every step that
replaces the witness, namely isolated exchange, the top shortening, column
shortening and supported compression, also preserves the positive condition.
Nothing below uses finiteness of an alphabet: block alphabets are infinite,
and only state sets, supports, type spaces and summary values must be finite.

\paragraph{The constants}
The bound $\Lensharp(N,d)$ of \cref{thm:mo-anchored} and the constants behind
it are as follows. For integers $N\ge1$, $d\ge0$, $M\ge1$ and $k\ge0$ put
\begin{equation}\label{eq:mo-constants}
\begin{aligned}
 \sCsharp(N,d)&=N^N\,2^{5N^2+10N(N+d)(d+1)+7N+1},&
 \sAsharp(N,d)&=d+1+\sCsharp(N,d),\\
 \Xisharp(N,d)&=6N^2+10N(N+d)(d+1)+7N+3d+4,&
 \sGsharp(N,d)&=2+\sAsharp(N,d)\,3^d,\\
 \Pcol(M,k)&=\mathsf{col}(M,k)^2\,2^{2M(M+1)k}\,3^k,\\
 \Lensharp(N,d)&=3(N+d+1)\,\Pcol\bigl(N+d+1,\sGsharp(N,d)\bigr)^d,
\end{aligned}
\end{equation}
with $\mathsf{col}$ from \cref{eq:bg-col}. They extend the constants
\cref{eq:bg-col} of the plain-entailment proof by the new tables and fields:
$\sCsharp(N,d)$ bounds the anchored exchange types at one level,
$\sAsharp(N,d)$ the protected and isolated letters of a normal hole at one
level, $\sGsharp(N,d)$ its dictionaries, $\Pcol(M,k)$ the positive column
summaries of $k$ blocks over a monitor with $M$ states, and $\Lensharp(N,d)$ is
the resulting length bound; $\Xisharp(N,d)$ is the exponent of its closed form.
The functions are nondecreasing in each argument, and the alphabet enters none
of them.

\subsection{The Separator Word}\label{app:mo-separator}

\begin{proof}[Proof of \cref{lem:mo-sep}]
(i) Let $y\le m$. The suffix $\wsuf wy$ contains $\sep$ at offset $m-y$. The
suffix $\wsuf wx$ is longer, and its letter at offset $m-y$ is $w[x+m-y]$ with
$x+m-y<m$, a letter of $\Sigma$. Hence $\wsuf wy\not\pre\wsuf wx$.

(ii) Here $\wsuf w{m+1+e}=\wsuf te$ is a word over $\Sigma$ and
$\wsuf wx=\wsuf ux\,\sep\,t$. If $\len{\wsuf te}\le\len{\wsuf ux}$, the prefix
test reads only letters of $\wsuf ux$, so it is the test
$\wsuf te\pre\wsuf ux$. Otherwise $\wsuf te\not\pre\wsuf ux$ by length, and also
$\wsuf te\not\pre\wsuf wx$, because the letter of $\wsuf wx$ at offset $m-x$ is
$\sep$, while the letter of $\wsuf te$ at that offset lies in $\Sigma$.

(iii) Here $\wsuf w{m+1+i}=\wsuf ti$ and $\wsuf w{m+1+j}=\wsuf tj$.

(iv) The letter $\sep$ occurs in $w$ only at $m$. The continuation letter of a
comparison $(x,m+1+e)$ is $w[x+\len w-m-1-e]=w[x+\len t-e]$, so the comparison
is anchored iff $x+\len t-e=m$, that is, iff $\len{\wsuf te}=\len{\wsuf ux}$.
By~(ii), the pair is a comparison iff $\wsuf te\pre\wsuf ux$, which together
with equal lengths says $\wsuf te=\wsuf ux$. Conversely, $\wsuf te=\wsuf ux$
gives equal lengths and, by~(ii), a comparison.
\end{proof}

We prove \cref{lem:mo-pi-reader} for the reader $\RPi$ and the relation $\Pi$
constructed in \cref{sec:mo-separator}. The cap $N_A$ of the counter is the
mortal bound of \cref{def:nm-mortal-hole}, not the size of $\RPi$; for $N_A=1$
no state $\sDone(r',k)$ is rejected, as required by $\len t-J(t)<1$.

\begin{proof}[Proof of \cref{lem:mo-pi-reader}]
The reader and $\Pi$ are computable, since mortality is decided by
reachability in the safe graph of $A$ (\cref{def:nm-kinds}).

\emph{Runs.} The states after the separator have no transition on $\sep$ other
than into $\sSink$, which is absorbing. So a word without $\sep$ ends in a
state $\sPre(\cdot)$ or in $\sSink$, and a word with two occurrences of $\sep$
ends in $\sSink$. Let $w=u\sep t$ with $u,t\in\Sigma^*$ and $m=\len u$. If $u$
is not safe, the run enters $\sSink$ within $u$ and stays there. If $u$ is
safe, then by induction the state at a cut $x\le m$ is
$\sPre(A_x(u),A'_x(u))$, and the separator leads to $\sAct$ of the start
states of $B$ and $B'$ iff $A(u)$ is mortal, and to $\sSink$ otherwise. In the
first case, by induction on $e$, the state at the cut $m+1+e$ is
$\sAct(B_e(t),B'_e(t))$ while no earlier cut of $t$ is rejected by $B$, that is,
for $e\le J(t)$ if $J(t)$ is defined and for all $e$ otherwise; and it is
$\sDone(B'_e(t),\min(e-J(t),N_A))$ for $e>J(t)$.

(i) There are $N_AN_{A'}$ states $\sPre(q,q')$, $N_BN_{B'}$ states
$\sAct(r,r')$, $N_{B'}N_A$ states $\sDone(r',k)$, and one sink.

(iii) and rejection. By the description of runs, the last state of a word is
rejected iff the word is $u\sep t$ with $u,t\in\Sigma^*$, $u$ is safe, $A(u)$
is mortal, $J(t)$ is defined, and $\len t-J(t)\le N_A-1$: if $\len t=J(t)$, the
run ends in $\sAct(B_{J(t)}(t),\cdot)$ with $B_{J(t)}(t)\in\Rej$; if $\len
t>J(t)$, it ends in a state $\sDone$ whose counter $\min(\len t-J(t),N_A)$ is
at most $N_A-1$ iff $\len t-J(t)\le N_A-1$. No state $\sPre$ is rejected, and
the states $\sAct(B_e(t),\cdot)$ with $e<J(t)$ are not, so the first rejected
cut is $m+1+J(t)$.

(ii) Let $w=u\sep t$ be rejected, as just described. By \cref{lem:mo-sep}, every
comparison of $w$ has the form $(x,m+1+e)$ with $x\le m$, a cross comparison
$(x,e)$ of $(u,t)$, or $(m+1+i,m+1+j)$ with $i<j$, a comparison $(i,j)$ of
$t$. Admitted pairs join a state $\sPre$ to a state $\sAct$, or two states
$\sAct$; the states $\sAct$ occur exactly at the cuts $m+1+e$ with $e\le J(t)$.
So a comparison $(x,m+1+e)$ is admitted iff $e\le J(t)$ and
$(A_x(u),B_e(t))\in\Guard$, that is, iff $(x,e)$ is an event of $(u,t)$; and a
comparison $(m+1+i,m+1+j)$ is admitted iff $j\le J(t)$ and
$(B_i(t),B_j(t))\in\GuardS$, that is, iff $(i,j)$ is a self event of $t$. Hence
$w$ is a hole iff $(u,t)$ has no event and $t$ has no self event. By
\cref{lem:mo-sep}(iv), the anchored comparisons of $w$ are the pairs
$(x,m+1+e)$ with $\wsuf ux=\wsuf te$. Their first state is
$\sPre(A_x(u),A'_x(u))$, and their second state is $\sAct(B_e(t),B'_e(t))$ or
$\sDone(B'_e(t),k)$, with $B'$-component $B'_e(t)$ in both cases, also when
$e>J(t)$. So $\Pi$ holds on such a pair iff $(A'_x(u),B'_e(t))\in\Head$, and an
anchored comparison in $\Pi$ is exactly a head event. Altogether, $w$ is a
$\Pi$-hole iff $w=u\sep t$ where $(u,t)$ satisfies conditions (a)--(e) of
\cref{def:ex-ss2self}, $A(u)$ is mortal and $\len t-J(t)<N_A$, that is, iff
$(u,t)$ is a mortal hole. A word that is not rejected is not a hole.
\end{proof}

\subsection{Positive Desubstitution}\label{app:mo-desub}

Throughout this subsection, $\mathcal M$ is a lettered monitor with a positive
channel $\Theta=(F,R,S,S',\Phi)$, $\ell$ is a marker, and
$v=X_0\ell X_1\ell\cdots X_{k-1}\ell$ is the expansion of a nonempty word
$v'=[X_0]\cdots[X_{k-1}]$ of blocks avoiding $\ell$. We use the notation of the
label facts \cref{eq:bg-label-facts}: the run $\zeta'_0=\star,\ldots,\zeta'_k$
of $v'$ in $\mathcal M_\ell$, the cuts $a_i=\len{X_0\ell\cdots X_{i-1}\ell}$ at
which the blocks begin, and $b_i=a_i+\len{X_i}$; in addition $a_k=\len v$. The
letter of $v$ at the position $a_i+g$, for $0\le g\le\len{X_i}$, is
$(X_i\ell)[g]$. We write $L$ and $L'$ for the labels of $v$ and $v'$.

\begin{lemma}[Transport of continuation letters]\label{lem:mo-transport}
Let $(x,y)$ be a comparison of $v$ with $y<\len v$, and let $i<i'<k$,
$X_i=\hat x\hat y$ and $X_{i'}=\hat x'\hat y$ be as in \cref{lem:bg-alignment},
so that $x=a_i+\len{\hat x}$, $y=a_{i'}+\len{\hat x'}$, and $(i+1,i'+1)$ is a
comparison of $v'$. Put $r=(i+1)+k-(i'+1)$, the position of the continuation
letter of $(i+1,i'+1)$ in $v'$. Then $i<r<k$, $x+\len v-y=a_r$, and the
continuation letter of $(x,y)$ is $v[a_r]=(X_r\ell)[0]$.
\end{lemma}

\begin{proof}
Since $\hat y$ ends both blocks, $b_i-x=\len{\hat y}=b_{i'}-y$. Hence
\[
 \len v-y=\len{\hat y}+1+\sum_{h=i'+1}^{k-1}(\len{X_h}+1)
 =\len{\hat y}+1+(a_k-a_{i'+1}).
\]
Since $(i+1,i'+1)$ is a comparison of $v'$, $X_{i+1+m}=X_{i'+1+m}$ for
$0\le m<k-i'-1$, so the sum equals $a_r-a_{i+1}$, where $r=i+k-i'$. With
$x+\len{\hat y}+1=b_i+1=a_{i+1}$ this gives $x+\len v-y=a_r$. Finally $i<r$
because $i'<k$, and $r<k$ because $i<i'$.
\end{proof}

For $i'+1=k$ the comparison $(i+1,k)$ of $v'$ is terminal and $r=i+1$: the
continuation letter is $v'[i+1]$, as in clause (S$'$).

\begin{proof}[Proof of \cref{lem:mo-desub}]
The second claim follows from the first and \cref{lem:bg-desub}. For the first,
recall from \cref{eq:bg-label-facts} that $L_{a_i+g}=\lab_g(\zeta'_i,X_i)$ for
$g\le\len{X_i}$, $L_{\len v}=(\zeta'_k,\ell)$, and $L'_{i+1}=(\zeta'_i,X_i)$,
with $\zeta'_k=\delta_\ell(\zeta'_{k-1},X_{k-1})$. Every cut $x<\len v$ lies in
exactly one interval $[a_i,b_i]$, and then $x=a_i+g$ with $g\le\len{X_i}$ and
$v[x]=(X_i\ell)[g]$.

($\Rightarrow$) Suppose $\PosOp_\Theta(v)$, and consider the clause that holds.
\begin{itemize}[leftmargin=*]
\item (F): $F(L_{\len v})=F(\delta_\ell(\zeta'_{k-1},X_{k-1}),\ell)$ is the first
  disjunct of $F_\ell(\zeta'_{k-1},X_{k-1})=F_\ell(L'_k)$, which is clause (F)
  of $v'$.
\item (R): a comparison $(x,y)$ with $y<\len v$, $R(L_x,L_y)$ and
  $v[x+\len v-y]\in\Phi$. Take $i<i'$, $\hat x$, $\hat x'$ and $\hat y$ as in
  \cref{lem:bg-alignment}. Then $L_x=\lab_{\len{\hat x}}(\zeta'_i,X_i)$ and
  $L_y=\lab_{\len{\hat x'}}(\zeta'_{i'},X_{i'})$, so $R(L_x,L_y)$ witnesses
  $R_\ell(L'_{i+1},L'_{i'+1})$. By \cref{lem:mo-transport}, the continuation
  letter of $(x,y)$ is $(X_r\ell)[0]$, so $X_r\in\Phi_\ell$, where $v'[r]$ is
  the continuation letter of the comparison $(i+1,i'+1)$ of $v'$. If
  $i'+1<k$, this is clause (R) of $v'$. If $i'+1=k$, then $r=i+1$, and since
  $S'_\ell=R_\ell$, this is clause (S$'$) of $v'$ at the cut $i+1<k$.
\item (S) or (S$'$): a cut $x<\len v$ with $S(L_x,L_{\len v})$, or with
  $S'(L_x,L_{\len v})$ and $v[x]\in\Phi$. Write $x=a_i+g$. Then
  $L_x=\lab_g(\zeta'_i,X_i)$, $L_{\len v}=(\zeta'_k,\ell)$ and
  $v[x]=(X_i\ell)[g]$, so $\ptest(\zeta'_i,X_i;\zeta'_k)$ holds. If $i<k-1$,
  this is $S_\ell((\zeta'_i,X_i),(\zeta'_{k-1},X_{k-1}))=S_\ell(L'_{i+1},L'_k)$,
  clause (S) of $v'$ at the cut $i+1<k$. If $i=k-1$, it is the second disjunct
  of $F_\ell(L'_k)$, clause (F) of $v'$.
\end{itemize}

($\Leftarrow$) Suppose $\PosOp_{\Theta_\ell}(v')$. All relations of $\Theta_\ell$
are false at $\triangleright=L'_0$, so a witnessing cut $x'$ of $v'$ satisfies
$x'\ge1$; write $x'=i+1$.
\begin{itemize}[leftmargin=*]
\item (F): $F_\ell(L'_k)$. Either $F(\zeta'_k,\ell)=F(L_{\len v})$, clause (F)
  of $v$; or $\ptest(\zeta'_{k-1},X_{k-1};\zeta'_k)$, which gives some
  $g\le\len{X_{k-1}}$ with $S(L_x,L_{\len v})$, or with $S'(L_x,L_{\len v})$ and
  $v[x]\in\Phi$, at the cut $x=a_{k-1}+g\le b_{k-1}<\len v$: clause (S) or
  (S$'$) of $v$.
\item (S): $S_\ell(L'_{i+1},L'_k)$ with $i+1<k$, that is,
  $\ptest(\zeta'_i,X_i;\zeta'_k)$. As in the previous case this gives clause (S)
  or (S$'$) of $v$ at a cut $x=a_i+g\le b_i<\len v$.
\item (R) or (S$'$): a comparison $(i+1,i'+1)$ of $v'$ with
  $R_\ell(L'_{i+1},L'_{i'+1})$ and continuation letter $v'[r]\in\Phi_\ell$,
  where either $i'+1<k$ and $r=i+k-i'$ (clause (R)), or $i'+1=k$, $r=i+1$, and
  the comparison is the terminal pair $(i+1,k)$, which is a comparison because
  $i+1<k$ (clause (S$'$), since $S'_\ell=R_\ell$). In both cases $i<i'<k$.
  The relation $R_\ell$ gives decompositions $X_i=\hat x\hat y$ and
  $X_{i'}=\hat x'\hat y$ with
  $R(\lab_{\len{\hat x}}(\zeta'_i,X_i),\lab_{\len{\hat x'}}(\zeta'_{i'},X_{i'}))$.
  Put $x=a_i+\len{\hat x}$ and $y=a_{i'}+\len{\hat x'}$. Then $y\le b_{i'}<\len
  v$, and by \cref{lem:bg-alignment} the pair $(x,y)$ is a comparison of $v$; by
  the label facts $R(L_x,L_y)$ holds; and by \cref{lem:mo-transport} its
  continuation letter is $(X_r\ell)[0]$, which lies in $\Phi$ because
  $X_r\in\Phi_\ell$. This is clause (R) of $v$.
\end{itemize}
This proves the equivalence.
\end{proof}

\subsection{Anchored Exchange and Bounded Dictionaries}\label{app:mo-exchange}

Fix a lettered monitor $\mathcal M_0$ with $N\ge1$ states and a positive
channel $\Theta_0$ on it, and a positive hole $w$ of depth $D\le d$ with
hierarchy $(w_i,\ell_i)$, monitors $\mathcal M_i$ and channels
$\Theta_i=(F_i,R_i,S_i,S'_i,\Phi_i)$ as in \cref{cor:mo-tower}. We use the
notions of step~2 of \cref{app:bg-search}: the dictionaries $\mathcal A_i$ and
$g_i=\len{\mathcal A_i}$, the start counts $s_i$ and start prefixes, the entry
states $U_i$ with $\len{U_i}\le N+s_i\le N+d$, the protected letters $P_i$ with
$\len{P_i}\le s_i+1\le d+1$, the boundary pairs $\operatorname{BP}_i$ with
$b_i=\len{\operatorname{BP}_i}$, and isolated letters. Recall that every
occurrence of an unprotected letter has an ordinary entry state.

\begin{definition}[Anchored exchange type]\label{def:mo-type}
Let $1\le i<D$. The \emph{anchored exchange type} of an unprotected letter $a$
of $w_i$ consists of its exchange type (\cref{app:bg-search}) together with
the following tables of $(\mathcal M_i,\Theta_i)$:
\begin{itemize}[leftmargin=*]
\item the row $F_i(\zeta,a)$ for every ordinary state $\zeta$;
\item for each $Y\in\{R_i,S_i,S'_i\}$, the self table $Y((\zeta,a),(\zeta',a))$
  for ordinary $\zeta,\zeta'$;
\item for each such $Y$, the protected rows $Y((\zeta,a),(\vartheta,z))$ and
  $Y((\vartheta,z),(\zeta,a))$ for ordinary $\zeta$, $\vartheta\in U_i$ and
  $z\in P_i$;
\item for each such $Y$, the initial row $Y(\triangleright,(\zeta,a))$ for
  ordinary $\zeta$, which is false at every level $i\ge1$;
\item the bit $[a\in\Phi_i]$.
\end{itemize}
\end{definition}

\begin{lemma}[Anchored exchange types]\label{lem:mo-types}
At every level $1\le i<D\le d$ there are at most $\sCsharp(N,d)$ anchored
exchange types.
\end{lemma}

\begin{proof}
Ordinary states transition to ordinary states, so the transition table has at
most $N^N$ values. The rows of $\Tfin_i$ and $F_i$ take $2N$ bits; the self
tables of the five relations $\Lambda_i,\Lfin_i,R_i,S_i,S'_i$ take $5N^2$ bits,
their protected rows $10N\len{U_i}\len{P_i}\le10N(N+d)(d+1)$ bits, and their
initial rows $5N$ bits; and the bit of $\Phi_i$ is one more. Hence there are at
most $N^N2^{5N^2+10N(N+d)(d+1)+7N+1}=\sCsharp(N,d)$ types.
\end{proof}

\begin{lemma}[Anchored isolated exchange]\label{lem:mo-exchange}
Let $1\le i<D$, let $a\ne b$ be isolated letters of $w_i$ with the same
anchored exchange type, where the level-zero expansion of $b$ is not longer
than that of $a$, let $h$ replace $a$ by $b$ and fix every other letter, and
let $w'$ be the expansion of $h(w_i)$ with the original lower markers. Then
$w'$ is a positive hole of $(\mathcal M_0,\Theta_0)$ of depth $D$, no longer
than $w$, with every start count unchanged, no dictionary larger than the
corresponding one of $w$, and level-$i$ dictionary $\mathcal A_i\setminus\{a\}$.
\end{lemma}

\begin{proof}
The anchored type contains the exchange type, so $a$ and $b$ have the same
exchange type, and by \cref{lem:bg-isolated-exchange} the map $h$ satisfies the
hypotheses of \cref{lem:bg-level-maps}: it identifies no unequal predecessor
pair of $v=w_i$, so $v$ and $h(v)$ have exactly the same comparisons
(\cref{lem:bg-comparison-map}); it fixes the start prefix; it preserves the
state run of $v$ in $\mathcal M_i$; and $h(v)$ is a hole of $\mathcal M_i$.
Hence $w'$ is a hole of $\mathcal M_0$ of depth $D$ with the stated start
counts, dictionaries and length, and its markers below $i$ are
$\ell_0,\ldots,\ell_{i-1}$, with $w'_i=h(v)$. So the monitors and channels of
$w'$ through level $i$ are $(\mathcal M_j,\Theta_j)_{j\le i}$, and by
\cref{lem:mo-desub}, applied at the levels $i-1,\ldots,0$, it remains to show
$\PosOp_{\Theta_i}(h(v))$.

Write $L$ and $L^h$ for the labels of $v$ and of $h(v)$ in $\mathcal M_i$. The
runs coincide, so $L^h_c$ is $L_c$ with its letter mapped by $h$.

\emph{Claim: for every comparison $(x,y)$ of $v$ and each
$Y\in\{R_i,S_i,S'_i\}$, $Y(L^h_x,L^h_y)\iff Y(L_x,L_y)$.} If $x=0$, then
$L_x=L^h_x=\triangleright$, and $L_y$ is unchanged unless $v[y-1]=a$, when the
initial row applies. Let $x\ge1$. If neither $v[x-1]$ nor $v[y-1]$ is $a$, both
labels are unchanged. If both are $a$, the self table applies, both entry
states being ordinary. If exactly one of them is $a$, the other letter $c$
differs from $a$, so $(x,y)$ is left maximal. If $y<\len v$, the ordered pair of
predecessor letters lies in $\operatorname{BP}_i$, and isolation of $a$ makes
$c$ protected. If $y=\len v$, then $v[y-1]=\ell_i$ is protected and differs
from $a$, so $v[x-1]=a$ and $c=\ell_i$. In both cases $c\in P_i$ is read from a
state in $U_i$, $c\ne b$ because $b$ is unprotected, and the protected row of
$Y$ in the appropriate orientation applies. This proves the claim.

Now fix a clause witnessing $\PosOp_{\Theta_i}(v)$.
\begin{itemize}[leftmargin=*]
\item (F): the final label $L_{\len v}$ has the protected letter $\ell_i$, so
  $L^h_{\len v}=L_{\len v}$.
\item (R): a comparison $(x,y)$ with $y<\len v$. It is a comparison of $h(v)$,
  the claim keeps $R_i(L_x,L_y)$, and the continuation letter is at the same
  position $x+\len v-y$, since $h$ preserves lengths. That letter is either
  unchanged or changed from $a$ to $b$, and $[a\in\Phi_i]=[b\in\Phi_i]$.
\item (S) or (S$'$): a cut $x<\len v$. The pair $(x,\len v)$ is a terminal
  comparison of both words, the claim keeps the value of $S_i$ or $S'_i$, and
  for (S$'$) the letter $h(v)[x]$ lies in $\Phi_i$ iff $v[x]$ does.
\end{itemize}
Hence $\PosOp_{\Theta_i}(h(v))$.
\end{proof}

The merge is not injective, so it gives no agreement of the derived data
above level $i$. After every merge we recompute the hierarchy, the towers, the
protected sets and the anchored types of the new word; its positivity in the
new towers follows from \cref{cor:mo-tower}.

A positive hole is \emph{anchored-normal} if at every level $1\le i<D$ distinct
isolated letters have distinct anchored exchange types, computed in its own
hierarchy and towers.

\begin{lemma}[Normal form and bounded dictionaries]\label{lem:mo-dictionary}
\leavevmode
\begin{enumerate}[leftmargin=*]
\item Every positive hole of depth $D$ has an anchored-normal positive hole of
  depth $D$, no greater length, the same start counts, and no larger dictionary
  at any level.
\item Let $w$ be an anchored-normal positive hole of depth $D\le d$, and put
  $\sAsharp=\sAsharp(N,d)$. Then $g_D=1$, and for $1\le i<D$,
  \[
   g_i\le\sAsharp+2b_i,\qquad
   2b_i\le(\sAsharp-1)(3^{D-1-i}-1),\qquad
   g_i\le1+(\sAsharp-1)3^{D-1-i}\le\sGsharp(N,d).
  \]
\end{enumerate}
Consequently, every positive hole of depth $D\le d$ has a positive hole of depth
$D$ and no greater length with $g_i\le\sGsharp(N,d)$ for every $1\le i\le D$.
\end{lemma}

\begin{proof}
(1) Use the potential $\sum_{j=0}^Dg_j$. While anchored normality fails at some
level $1\le i<D$, merge an offending pair by \cref{lem:mo-exchange}, in the
direction of the shorter level-zero expansion. The result is a positive hole of
depth $D$, no longer, with the same start counts; no dictionary grows and the
level-$i$ dictionary shrinks, so the potential decreases strictly. The process
terminates in an anchored-normal positive hole.

(2) The top level $w_D$ is unary, so $g_D=1$, and $b_D=0$ because a unary word
has no left-maximal comparison. Let $1\le i<D$. Partition $\mathcal A_i$ into
the protected letters, at most $d+1$; the isolated letters, at most one per
anchored type by normality, hence at most $\sCsharp(N,d)$ by
\cref{lem:mo-types}; and the remaining letters, each of which is unprotected and
not isolated, hence a component of a pair in $\operatorname{BP}_i$. The two
projections of $\operatorname{BP}_i$ have at most $2b_i$ elements, so
$g_i\le\sAsharp+2b_i$. By the propagation in \cref{lem:bg-dictionary}(2), which
holds for every nonempty word, $b_j\le b_{j+1}+g_{j+1}-1$ for every $j$. Hence
$b_{D-1}\le b_D+g_D-1=0$, and for $1\le i<D-1$,
$b_i\le b_{i+1}+\sAsharp+2b_{i+1}-1=3b_{i+1}+\sAsharp-1$. By downward induction
on $i$, with $r=D-1-i$,
\[
 2b_i\le3(\sAsharp-1)(3^{r-1}-1)+2(\sAsharp-1)=(\sAsharp-1)(3^r-1),
\]
the base case $r=0$ having both sides zero. Then
$g_i\le\sAsharp+2b_i\le1+(\sAsharp-1)3^r\le2+\sAsharp3^d=\sGsharp(N,d)$.

The last assertion follows from~(1) and~(2), since $g_D=1\le\sGsharp(N,d)$.
\end{proof}

\subsection{The Unary Top Level}\label{app:mo-unary}

\begin{lemma}[Positive unary words]\label{lem:mo-unary}
Let $\mathcal M$ be a lettered monitor with state set $\mathcal Q$, $\Theta$ a
positive channel on it, and $z^m$ a positive hole of $(\mathcal M,\Theta)$.
Then $z^{m'}$ is a positive hole of $(\mathcal M,\Theta)$ for some
$1\le m'\le3\len{\mathcal Q}$.
\end{lemma}

\begin{proof}
Every pair of cuts $x<y$ of a unary word is a comparison, and its continuation
letter is $z$. Fix a clause witnessing $\PosOp_\Theta(z^m)$, and let $W$ be the
set consisting of $m$ and the cuts of the witness other than $0$: the cuts $x$
and $y$ for (R), the cut $x$ for (S) and (S$'$), and none for (F); so $W\cup\{0\}$
has at most four elements.

Let $i<i'\le m-1$ be cuts with equal states such that no cut of $W$ lies in
$\{i+1,\ldots,i'\}$, and delete the letters at the positions $i,\ldots,i'-1$.
The cuts of the resulting word $z^{m-(i'-i)}$ correspond to the cuts $c\le i$
and $c>i'$ of $z^m$, by $c\mapsto c$ and $c\mapsto c-(i'-i)$, an
order-preserving map that keeps the labels: the state at the new cut $i$ is the
state at the old cut $i'$. The last cut maps to the last cut. Every pair of
cuts of a unary word is a comparison, so every comparison of the new word has
the labels of a comparison of $z^m$, terminal iff that one is, and the final
label is unchanged; hence the new word is a hole. The cuts of $W$ are kept,
with their labels and their order, and since every pair of cuts is a
comparison with continuation letter $z$, the witnessing clause holds for the
images of its cuts. So the shorter word is again a positive hole, and $W$ maps
onto the corresponding set of the new witness.

Repeat such deletions while possible; the length decreases, so this terminates,
in a positive hole $z^{m'}$ with $m'\ge1$, because the last letter is never
deleted. Let $c<c'$ be consecutive elements of $W\cup\{0\}$. The states at the
cuts $c,\ldots,c'-1$ are pairwise distinct: two equal ones at $i<i'$ would have
$i'\le c'-1\le m'-1$ and no cut of $W$ in $\{i+1,\ldots,i'\}$, so a further
deletion would be possible. Hence $c'-c\le\len{\mathcal Q}$, and, as there are at
most three such gaps and they cover $[0,m']$, $m'\le3\len{\mathcal Q}$.
\end{proof}

\subsection{Agreement with Channels}\label{app:mo-agreement}

We extend agreement \cref{eq:bg-agreement} to channels. Let $\mathcal U$ and
$\mathcal V$ be lettered monitors with positive channels
$\Theta^{\mathcal U}=(F^{\mathcal U},R^{\mathcal U},S^{\mathcal U},
S'^{\mathcal U},\Phi^{\mathcal U})$ and $\Theta^{\mathcal V}$, let $E$ be a finite
set of letters of $\mathcal U$, and $\rho$ a map from $E$ into the alphabet of
$\mathcal V$. We write
$(\mathcal U,\Theta^{\mathcal U})\equiv_{E,\rho}(\mathcal V,\Theta^{\mathcal V})$
if $\mathcal U\equiv_{E,\rho}\mathcal V$ and, for every $a\in E$, every state
$\zeta$, and all labels $L,L'$ whose letters lie in $E$, the start label
included,
\[
 F^{\mathcal V}(\zeta,\rho(a))\iff F^{\mathcal U}(\zeta,a),\qquad
 \rho(a)\in\Phi^{\mathcal V}\iff a\in\Phi^{\mathcal U},
\]
and $Y^{\mathcal V}(\widehat\rho L,\widehat\rho L')\iff Y^{\mathcal U}(L,L')$ for
$Y\in\{R,S,S'\}$.

\emph{Agreement preserves positive holes.} Let $v$ be a nonempty word over $E$.
Agreement of the monitors gives the same states along $v$ and $\rho(v)$, labels
that correspond under $\widehat\rho$, and holes that correspond. Since $\rho$ is
injective on $E$, the words $v$ and $\rho(v)$ have the same letter equalities,
hence the same comparisons, and the continuation letter of a comparison of
$\rho(v)$ is the image of that of $v$. Each clause of the positive condition
therefore holds for $v$ iff it holds for $\rho(v)$, and $v$ is a positive hole
of $(\mathcal U,\Theta^{\mathcal U})$ iff $\rho(v)$ is a positive hole of
$(\mathcal V,\Theta^{\mathcal V})$. \emph{Agreements compose}, by composing the
maps.

\begin{lemma}[Agreement under derivation]\label{lem:mo-agreement}
Suppose $(\mathcal U,\Theta^{\mathcal U})\equiv_{E,\rho}(\mathcal
V,\Theta^{\mathcal V})$ and $\ell\in E$. For every finite family $\mathcal F$ of
blocks over $E$ avoiding $\ell$, the letterwise map $\rho^*$ gives
\[
 (\mathcal U_\ell,\Theta^{\mathcal U}_\ell)\equiv_{\mathcal F,\rho^*}
 (\mathcal V_{\rho(\ell)},\Theta^{\mathcal V}_{\rho(\ell)}),
\]
and the image blocks avoid $\rho(\ell)$.
\end{lemma}

\begin{proof}
The agreement of the derived monitors and the avoidance of $\rho(\ell)$ are
\cref{lem:bg-agreement}. Its proof shows, by induction along a block, that for
$X\in\mathcal F$ and $\zeta\in\mathcal Q\sqcup\{\star\}$ the shifted blocks of
$X$ and $\rho^*(X)$ have the same states, so that
$\lab_g(\zeta,\rho^*X)=\widehat\rho\,\lab_g(\zeta,X)$ for every $g$ and the exits
coincide. The letters $(X\ell)[g]$ lie in $E$, and
$(\rho^*(X)\rho(\ell))[g]=\rho((X\ell)[g])$. Hence:
\begin{itemize}[leftmargin=*]
\item $\rho^*(X)\in\Phi^{\mathcal V}_{\rho(\ell)}$ iff
  $\rho((X\ell)[0])\in\Phi^{\mathcal V}$ iff $(X\ell)[0]\in\Phi^{\mathcal U}$
  iff $X\in\Phi^{\mathcal U}_\ell$;
\item by injectivity, the decompositions $X=xy$, $X'=x'y$ correspond exactly to
  the decompositions of $\rho^*X$ and $\rho^*X'$ with a common suffix, and
  $R$ agrees at the corresponding labels, whose letters lie in $E$ or which are
  $\triangleright$; so $R_\ell$, and hence $S'_\ell$, agree;
\item the positive terminal test $\ptest(\zeta,X;\zeta'')$ reads $S$ and $S'$
  at the labels $\lab_g(\zeta,X)$ and $(\zeta'',\ell)$, which correspond to
  $\lab_g(\zeta,\rho^*X)$ and $(\zeta'',\rho(\ell))$, and the letters
  $(X\ell)[g]$; so it agrees, and with equal exits so do $S_\ell$ and, using
  agreement of $F$ at the letter $\ell\in E$, $F_\ell$.
\end{itemize}
All relations at the new start label are false on both sides.
\end{proof}

\subsection{Positive Column Summaries}\label{app:mo-columns}

Fix a lettered monitor $\mathcal M$ with state set $\mathcal Q$,
$M=\len{\mathcal Q}$, start $\zeta_0$, a positive channel
$\Theta=(F,R,S,S',\Phi)$, a marker $\ell$, and blocks $X_1,\ldots,X_k$ avoiding
$\ell$. We use the summary, the column word and the notation $\lab_0(\zeta)$ of
step~4 of \cref{app:bg-search}; block indices are $a,b$, entry states
$\zeta,\zeta'$ range over $\mathcal Q\sqcup\{\star\}$, and final states $\zeta''$
over $\mathcal Q$.

\paragraph{Fields}
Let $\mathcal M^\Theta$ be the lettered monitor with the states, start and
transitions of $\mathcal M$, the final target $F$, and the relations $R$ and $S$
in place of $\Lambda$ and $\Lfin$. The \emph{positive summary} of
$(X_1,\ldots,X_k)$ consists of the following fields:
\begin{center}
\small
\begin{tabular}{@{}lp{.7\linewidth}@{}}
\toprule
Field & Meaning\\
\midrule
$\mathit{sum}$ & the summary of \cref{app:bg-search} of $(X_1,\ldots,X_k)$ for $\mathcal M$\\
$\mathit{sum}^\Theta$ & the summary of \cref{app:bg-search} of $(X_1,\ldots,X_k)$ for $\mathcal M^\Theta$\\
$\ftSp_a(\zeta,\zeta'')$ & some $g<\len{X_a}$ has
  $S'(\lab_g(\zeta,X_a),(\zeta'',\ell))$ and $X_a[g]\in\Phi$\\
$\fpend_a(\zeta,\zeta'')$ & $S'(\lab_{\len{X_a}}(\zeta,X_a),(\zeta'',\ell))$\\
$\ffst_a$ & $\mathit{none}$ if $X_a=\eps$, and otherwise the bit $[X_a[0]\in\Phi]$\\
\bottomrule
\end{tabular}
\end{center}
In $\mathit{sum}^\Theta$, the field $\fterm_a(\zeta,\zeta'')$ says that some
$g\le\len{X_a}$ has $S(\lab_g(\zeta,X_a),(\zeta'',\ell))$, and the field
$\fadm(a,b,\zeta,\zeta')$ says that some $X_a=xy$, $X_b=x'y$ has
$R(\lab_{\len x}(\zeta,X_a),\lab_{\len{x'}}(\zeta',X_b))$; its exits and
equality bits repeat those of $\mathit{sum}$. The field $\fpend$ records an
$S'$-test at the last position of the shifted block, whose letter is not known
yet: it is the next letter of the block, or the marker if the block ends there.
Counting assignments, the positive summary takes at most $\Pcol(M,k)$ values:
each of the two summaries takes at most $\mathsf{col}(M,k)$ values
(\cref{app:bg-constants}); the fields $\ftSp$ and $\fpend$ give
$2^{2M(M+1)k}$, since $a$ ranges over $k$ blocks, $\zeta$ over $M+1$ entry
states and $\zeta''$ over $M$ states; and $\ffst$ gives $3^k$.

\paragraph{The automaton}
The tuple of empty blocks has the initial summaries of \cref{app:bg-search}
for $\mathcal M$ and for $\mathcal M^\Theta$, $\ftSp_a$ false,
$\fpend_a(\zeta,\zeta'')=S'(\lab_0(\zeta),(\zeta'',\ell))$, and
$\ffst_a=\mathit{none}$. Reading a column $\gamma$ updates both summaries as in
\cref{app:bg-search}, and the other fields as follows, simultaneously and from
the old values. If $\gamma_a$ is absent, the fields of block $a$ are unchanged.
If $\gamma_a=z$, then, with $e=\fexit_a(\zeta)$ the old exit,
\[
\begin{aligned}
 \ftSp'_a(\zeta,\zeta'')&=\ftSp_a(\zeta,\zeta'')\lor
   \bigl(\fpend_a(\zeta,\zeta'')\land z\in\Phi\bigr),\qquad
 \fpend'_a(\zeta,\zeta'')=S'((e,z),(\zeta'',\ell)),\\
 \ffst'_a&=\ffst_a\text{ if }\ffst_a\ne\mathit{none},\ \text{and }[z\in\Phi]\text{ otherwise}.
\end{aligned}
\]

\begin{lemma}[Positive summaries]\label{lem:mo-summary}
\leavevmode
\begin{enumerate}[leftmargin=*]
\item After every prefix of a valid column word, every field equals its meaning
  for the current tuple.
\item The final positive summary of $(X_1,\ldots,X_k)$ determines the derived
  channel $\Theta_\ell$ on the letters $X_1,\ldots,X_k$, at every state of
  $\mathcal M_\ell$. With $\fterm^\Theta$ and $\fadm^\Theta$ the fields of
  $\mathit{sum}^\Theta$ and
  \[
   T_a(\zeta,\zeta'')=\fterm^\Theta_a(\zeta,\zeta'')\lor\ftSp_a(\zeta,\zeta'')\lor
   \bigl(\fpend_a(\zeta,\zeta'')\land\ell\in\Phi\bigr),
  \]
  it is given by
  \[
  \begin{aligned}
   X_a\in\Phi_\ell&\iff\ffst_a=1\ \lor\ (\ffst_a=\mathit{none}\land\ell\in\Phi),\\
   R_\ell((\zeta,X_a),(\zeta',X_b))\iff S'_\ell((\zeta,X_a),(\zeta',X_b))
     &\iff\fadm^\Theta(a,b,\zeta,\zeta'),\\
   S_\ell((\zeta,X_a),(\zeta',X_b))&\iff T_a(\zeta,\fexit_b(\zeta')),\\
   F_\ell(\zeta,X_a)&\iff F(\fexit_a(\zeta),\ell)\lor
     T_a(\zeta,\fexit_a(\zeta)).
  \end{aligned}
  \]
  In particular $S_\ell$ has no $F$-term. All relations at the start label of
  $\mathcal M_\ell$ are false by definition.
\end{enumerate}
\end{lemma}

\begin{proof}
(1) Induction on the number of columns read. For the two summaries this is
the correctness of the column automaton of \cref{app:bg-search}, applied to
$\mathcal M$ and to $\mathcal M^\Theta$. The initial values of the other fields
are their meanings for empty blocks: the shifted block of $\eps$ read from
$\zeta$ has the single label $\lab_0(\zeta)$, and no position $g<0$ exists. An
absent entry leaves the prefix of its track unchanged. Let track $a$ read $z$,
so that its prefix $X$ becomes $Xz$. The labels of the shifted block of $Xz$
are those of $X$ at the positions $g\le\len X$, and the new last label
$\lab_{\len X+1}(\zeta,Xz)=(\fexit_a(\zeta),z)$. The letter at position $\len X$
of $Xz$ is $z$, and the letters at the positions $g<\len X$ are unchanged.
This proves the updates of $\ftSp$ and $\fpend$: the new position $\len X$ of
$\ftSp$ carries the old pending test $\fpend$ and the letter $z$. By validity,
letters are appended to a track in order, so the first letter of a nonempty
track is the letter read when $\ffst$ leaves $\mathit{none}$.

(2) By~(1), the final fields are their meanings for $(X_1,\ldots,X_k)$. In
\cref{def:mo-channel-derivative}, $(X_a\ell)[0]$ is $X_a[0]$ if $X_a\ne\eps$ and
$\ell$ otherwise, which gives $\Phi_\ell$. The field $\fadm^\Theta$ is the
defining formula of $R_\ell=S'_\ell$. The positive terminal test is
$\ptest(\zeta,X_a;\zeta'')=T_a(\zeta,\zeta'')$, because $(X_a\ell)[g]$ is
$X_a[g]$ for $g<\len{X_a}$ and $\ell$ for $g=\len{X_a}$. Finally
$\delta_\ell(\zeta',X_b)=\fexit_b(\zeta')$.
\end{proof}

The summary of $\mathcal M$ alone does not suffice, as \cref{ex:mo-loses} shows.
We now give the field $\fpend$ in that example.

\begin{example}[\cref{ex:mo-loses} continued: the pending test]\label{ex:mo-loses-fields}
In \cref{ex:mo-loses} the top letter $[ab]$ is shortened as the single track
$ab$ over the level-$1$ monitor $\mathcal M_1$ and channel $\Theta_1$, with the
marker $c=[\sep1]$. For every state $\zeta''$ of $\mathcal M_1$, the shifted block
of $c$ read from $\zeta''$ reads $0\sep1$ from an ordinary state, or $\sep1$
from the start, and in both cases the state before the final $1$ is $4$. After
the prefix $a$, the last label of the shifted track is $(2,a)$ from an ordinary
entry, or $(\star_1,a)$ from the fresh entry, and the only common suffix of the
blocks $\eps$ and $\sep1$ is empty. So $\fpend$, the value of
$S'_1=R_1$ at that label and the final label $(\zeta'',c)$, compares the states
$\delta(2,0)=0$, or the start state $0$, with $\delta(4,1)=2$; as
$(0,2)\notin\Pi$, $\fpend$ is false everywhere. After the prefix $ab$ the last
label is $(0,b)$, and the common suffix $1$ of the blocks $1$ and $\sep1$
compares the states $\delta(0,0)=3$ and $4$; as $(3,4)\in\Pi$, $\fpend$ is true
everywhere. This is the original witness $(1,4)$, and it enters $F_2$ through
the disjunct $\fpend\land c\in\Phi_1$ of \cref{lem:mo-summary}(2), since the
expansion of $c$ starts with $\sep$. Hence the positive summaries after $a$ and
after $ab$ differ, and the column holding $b$ is kept. At level $0$ the tracks
$\eps$, $1$ and $\sep1$ admit no deletion either: their column word has two
columns, each of which is the first column of a nonempty track, where its field
$\ffst$ changes. The compression returns $w$ itself.
\end{example}

\begin{lemma}[Positive column shortening]\label{lem:mo-columns}
Let $\mathcal F=\{X_1,\ldots,X_k\}$ be a set of $k$ distinct blocks avoiding
$\ell$. There are subsequences $\tau(X_a)$ of the blocks $X_a$ with
$\len{\tau(X_a)}<\Pcol(M,k)$ for every $a$, such that $\tau$ is injective on
$\mathcal F$, the blocks $\tau(X_a)$ avoid $\ell$, and
\[
 (\mathcal M_\ell,\Theta_\ell)\equiv_{\mathcal F,\tau}(\mathcal M_\ell,\Theta_\ell).
\]
\end{lemma}

\begin{proof}
The new positive summary is a function of the old one and the column, and there
are at most $\Pcol(M,k)$ positive summaries. If the column word has at least
$\Pcol(M,k)$ columns, two of its cuts carry the same positive summary. Deleting
the columns between them leaves a valid column word, and the same final
summary, because the remaining columns are read from the same summary.
Repeating, we reach fewer than $\Pcol(M,k)$ columns; the tracks are
subsequences of the blocks and avoid $\ell$. Equal final summaries give equal
equality bits, so $\tau$ preserves equality and inequality of blocks and is
injective on $\mathcal F$; by \cref{lem:bg-columns} and \cref{app:bg-search},
the summary of $\mathcal M$ determines $\mathcal M_\ell$ on the letters, and by
\cref{lem:mo-summary}(2) the positive summary determines $\Theta_\ell$ on them,
at every state and at labels with letters in $\mathcal F$ or the start label.
Both sides are the same monitor and channel, with the same states and start.
\end{proof}

\subsection{Supported Compression}\label{app:mo-compression}

For markers $\boldsymbol\ell=(\ell_0,\ldots,\ell_{j-1})$ write
$(\mathcal M^{\boldsymbol\ell}_0,\Theta^{\boldsymbol\ell}_0)=(\mathcal
M_0,\Theta_0)$ and $(\mathcal M^{\boldsymbol\ell}_{h+1},
\Theta^{\boldsymbol\ell}_{h+1})=((\mathcal M^{\boldsymbol\ell}_h)_{\ell_h},
(\Theta^{\boldsymbol\ell}_h)_{\ell_h})$, and use valid families and the
expansions $\Exp_{\boldsymbol\ell,j}$ as in \cref{app:bg-search}.

\begin{lemma}[Supported compression with channels]\label{lem:mo-compression}
Fix a nonempty word $w$ with hierarchy markers $\boldsymbol\ell$ and dictionaries
$\mathcal A_h$, and $j\le d$. Suppose $\len{\mathcal A_h}\le g$ for $1\le h\le
j$, and put $\kappa_\sharp=\Pcol(N+d+1,g)$. For every set $E_j\subseteq
\mathcal A_j$ there are new markers $\widetilde{\boldsymbol\ell}=
(\widetilde\ell_0,\ldots,\widetilde\ell_{j-1})$ and a map $\rho_j$, injective on
$E_j$, such that
\[
 (\mathcal M^{\boldsymbol\ell}_j,\Theta^{\boldsymbol\ell}_j)\equiv_{E_j,\rho_j}
 (\mathcal M^{\widetilde{\boldsymbol\ell}}_j,\Theta^{\widetilde{\boldsymbol\ell}}_j),
 \qquad
 \len{\Exp_{\widetilde{\boldsymbol\ell},j}(\rho_j(a))}\le\kappa_\sharp^j
 \quad(a\in E_j),
\]
and the image family is valid for the new markers.
\end{lemma}

\begin{proof}
This is the induction of \cref{lem:bg-compression}, with positive
summaries and agreement with channels. For $j=0$ take the identity: a letter
expands to itself, of length $1=\kappa_\sharp^0$, whatever the size of $E_0$.

For the step from $j$ to $j+1\le d$, let $E_{j+1}\subseteq\mathcal A_{j+1}$, so
$\len{E_{j+1}}\le g$. Put $(\mathcal U,\Theta^{\mathcal U})=(\mathcal
M^{\boldsymbol\ell}_j,\Theta^{\boldsymbol\ell}_j)$, which has $N+j$ states. The
blocks in $E_{j+1}$ are distinct words over $\mathcal A_j$ avoiding $\ell_j$.
\Cref{lem:mo-columns} gives subsequences $\tau(X)$ with
$\len{\tau(X)}<\Pcol(N+j,\len{E_{j+1}})\le\kappa_\sharp$, by monotonicity,
and $(\mathcal U_{\ell_j},\Theta^{\mathcal U}_{\ell_j})\equiv_{E_{j+1},\tau}
(\mathcal U_{\ell_j},\Theta^{\mathcal U}_{\ell_j})$. Put
\[
 E_j=\{\ell_j\}\cup\bigcup_{X\in E_{j+1}}\alphof(\tau(X))\ \subseteq\ \mathcal A_j .
\]
The induction hypothesis gives markers $\widetilde\ell_0,\ldots,
\widetilde\ell_{j-1}$, the monitor and channel $(\mathcal V,\Theta^{\mathcal
V})=(\mathcal M^{\widetilde{\boldsymbol\ell}}_j,
\Theta^{\widetilde{\boldsymbol\ell}}_j)$, and $\rho_j$ with $(\mathcal
U,\Theta^{\mathcal U})\equiv_{E_j,\rho_j}(\mathcal V,\Theta^{\mathcal V})$ and
expansions of length at most $\kappa_\sharp^j$. Set
$\widetilde\ell_j=\rho_j(\ell_j)$ and $\rho_{j+1}(X)=\rho_j^*(\tau(X))$ for
$X\in E_{j+1}$. By \cref{lem:mo-agreement}, applied to the family
$\tau(E_{j+1})$ of blocks over $E_j$ avoiding $\ell_j$,
$(\mathcal U_{\ell_j},\Theta^{\mathcal U}_{\ell_j})\equiv_{\tau(E_{j+1}),
\rho_j^*}(\mathcal V_{\widetilde\ell_j},\Theta^{\mathcal V}_{\widetilde\ell_j})$,
and composition with the column agreement gives
$(\mathcal M^{\boldsymbol\ell}_{j+1},\Theta^{\boldsymbol\ell}_{j+1})
\equiv_{E_{j+1},\rho_{j+1}}(\mathcal M^{\widetilde{\boldsymbol\ell}}_{j+1},
\Theta^{\widetilde{\boldsymbol\ell}}_{j+1})$. The image blocks avoid
$\widetilde\ell_j$ and their letters lie in the valid family $\rho_j(E_j)$, so
the new family is valid. Each new block followed by its marker has at most
$\kappa_\sharp$ letters, each expanding to at most $\kappa_\sharp^j$ letters, so
every new letter expands to at most $\kappa_\sharp^{j+1}$ letters. No dictionary
or normality property of an intermediate compressed word is used.
\end{proof}

\subsection{Proof of the Anchored Compression Theorem}\label{app:mo-proof}

Let $\mathcal M_0$ be a lettered monitor with $N\ge1$ states and $\Theta_0$ a
positive channel on it, and suppose that $(\mathcal M_0,\Theta_0)$ has a
positive hole of depth at most $d$. By \cref{lem:mo-dictionary}, it has a
positive hole $w$ of some depth $D\le d$ with $g_i\le\sGsharp(N,d)$ for $1\le
i\le D$. Keep $w$, its hierarchy $(w_i,\ell_i)$ and its dictionaries fixed, and
put $\kappa_\sharp=\Pcol(N+d+1,\sGsharp(N,d))\ge1$.

\emph{The top level.} By \cref{cor:mo-tower}, the unary word $w_D=z^m$ is a
positive hole of $(\mathcal M_D,\Theta_D)$, and $\mathcal M_D$ has $N+D$ states.
By \cref{lem:mo-unary}, $z^{m'}$ is a positive hole of $(\mathcal
M_D,\Theta_D)$ for some $1\le m'\le3(N+D)$.

\emph{Compression.} Apply \cref{lem:mo-compression} with $j=D$,
$g=\sGsharp(N,d)$ and $E_D=\{z\}$. It gives markers $\widetilde{\boldsymbol\ell}$
and $\rho_D$ with $(\mathcal M_D,\Theta_D)\equiv_{\{z\},\rho_D}(\mathcal
M^{\widetilde{\boldsymbol\ell}}_D,\Theta^{\widetilde{\boldsymbol\ell}}_D)$.
Agreement preserves positive holes, so $\widetilde w_D=\rho_D(z)^{m'}$ is a
positive hole of $(\mathcal M^{\widetilde{\boldsymbol\ell}}_D,
\Theta^{\widetilde{\boldsymbol\ell}}_D)$.

\emph{Expansion.} For $i<D$ let $\widetilde w_i$ be the expansion of
$\widetilde w_{i+1}$ with the marker $\widetilde\ell_i$. By validity,
$\widetilde w_{i+1}$ is a nonempty word of blocks avoiding $\widetilde\ell_i$, so
\cref{lem:mo-desub} shows, by downward induction, that $\widetilde w_i$ is a
positive hole of $(\mathcal M^{\widetilde{\boldsymbol\ell}}_i,
\Theta^{\widetilde{\boldsymbol\ell}}_i)$. At $i=0$, $\widetilde w=\widetilde w_0$
is a positive hole of $(\mathcal M_0,\Theta_0)$. Its length is
\[
 m'\,\len{\Exp_{\widetilde{\boldsymbol\ell},D}(\rho_D(z))}
 \le3(N+D)\,\kappa_\sharp^D\le3(N+d+1)\,\kappa_\sharp^d=\Lensharp(N,d).
\]

\emph{Depth.} Each $\widetilde w_i$ with $i<D$ ends with its marker
$\widetilde\ell_i$, and cutting it at that marker returns $\widetilde w_{i+1}$.
So the canonical hierarchy of $\widetilde w$ is $\widetilde w_0,\ldots,
\widetilde w_D$ through level $D$, and $\widetilde w_D$ is unary. Hence
$\depth(\widetilde w)\le D\le d$.

If $D=0$, there is no compression: $w$ is unary and \cref{lem:mo-unary} gives a
positive hole of length at most $3N\le\Lensharp(N,0)$ and depth $0$.

This proves the general statement of \cref{thm:mo-anchored}. For a reader with
a relation $\Pi$, apply it to the lettered monitor of the reader and the base
channel $\Theta_\Pi$, whose positive holes are the $\Pi$-holes
(\cref{lem:mo-base}); the monitor has the $N$ states of the reader, and a
nonempty word has the same depth whichever monitor reads it. The estimates are
\cref{lem:mo-closed-form} below.

\subsection{Estimates and Decision Procedures}\label{app:mo-estimates}

\begin{lemma}[Closed form]\label{lem:mo-closed-form}
For $N\ge1$ and $d\ge0$, with $\Xisharp=\Xisharp(N,d)$,
\[
 \sGsharp(N,d)\le2^{\Xisharp}
 \qquad\text{and}\qquad
 \Lensharp(N,d)\le2^{2^{5\Xisharp+6}} .
\]
\end{lemma}

\begin{proof}
Put $\Xi_0=\Xisharp-3d-3=6N^2+10N(N+d)(d+1)+7N+1$. Since $N\le2^N$, we have
$N^N\le2^{N^2}$, so $\sCsharp(N,d)\le2^{\Xi_0}$. For nonnegative $x$ and $y$,
$2^x+2^y\le2^{x+y+1}$; with $d+1\le2^{d+1}$ this gives
$\sAsharp(N,d)\le2^{\Xi_0+d+2}$. Using $3^d\le2^{2d}$,
\[
 \sGsharp(N,d)=2+\sAsharp(N,d)\,3^d\le2+2^{\Xi_0+3d+2}\le2^{\Xi_0+3d+3}
 =2^{\Xisharp}.
\]
Put $H=\Xisharp$, $M=N+d+1$, $g=\sGsharp(N,d)$ and $P=2^H$. Then $M+1\le H$,
so $M$, $M+1$ and $d$ are at most $P$, and $g\le P$. Replacing $M$ by $2^M$ in
the first factor of \cref{eq:bg-col} gives
$\mathsf{col}(M,g)\le2^{2M(M+1)g+g^2+(M+1)^2g^2}\le2^{4P^4}=2^{2^{4H+2}}$.
Moreover $2^{M(M+1)g}\le2^{P^3}\le2^{2^{4H+2}}$ and
$3^g\le2^{2g}\le2^{2^{4H+2}}$. Hence the five factors of
$\Pcol(M,g)=\mathsf{col}(M,g)^2\,(2^{M(M+1)g})^2\,3^g$ give
$\Pcol(M,g)\le2^{5\cdot2^{4H+2}}\le2^{2^{4H+5}}$. Finally $3M\le2^{H+2}$ and
$d\le2^H$, so
\[
 \Lensharp(N,d)=3M\,\Pcol(M,g)^d\le2^{H+2+d\,2^{4H+5}}\le
 2^{2^{5H+5}+2^{5H+5}}=2^{2^{5H+6}}. \qedhere
\]
\end{proof}

\begin{proof}[Proof of \cref{lem:mo-estimate}]
Put $N_0=2^\alpha$ and $D_0=2^{\alpha+\beta+1}$. Then $\Nbar\le N_0$ and
$d_I=2\Nbar N_A-1<D_0$, so, by monotonicity and \cref{lem:mo-closed-form},
$\Lensharp_I\le\Lensharp(N_0,D_0)\le2^{2^{5\Xisharp(N_0,D_0)+6}}$. Since
$\alpha\ge3$ and $\beta\ge0$, we have $N_0\ge8$, $D_0\ge16$ and $N_0\le D_0/2$.
Hence $10N_0(N_0+D_0)(D_0+1)\le10N_0\cdot\frac32D_0\cdot\frac{17}{16}D_0<16N_0D_0^2$,
$6N_0^2\le3N_0D_0\le\frac3{16}N_0D_0^2$, and
$7N_0+3D_0+4\le N_0D_0\le\frac1{16}N_0D_0^2$, so that
\[
 \Xisharp(N_0,D_0)\le17N_0D_0^2=17\cdot2^{3\alpha+2\beta+2}
 \qquad\text{and}\qquad
 5\Xisharp(N_0,D_0)+6\le340\cdot2^{3\alpha+2\beta}+6<2^{3\alpha+2\beta+9}.
\]
Therefore $\Lensharp_I\le2^{2^{2^{3\alpha+2\beta+9}}}=E_3(3\alpha+2\beta+9)$.
\end{proof}

\begin{proof}[Proof of \cref{cor:mo-decidable}]
Compute $\Lensharp_I$, and test every pair $(u,t)$ of words over $\Sigma$ with
$\len u+1+\len t\le\Lensharp_I$: check conditions (a)--(e) of \cref{def:ex-ss2self}
over the finitely many cuts, decide whether $A(u)$ is mortal by reachability in
the safe graph, and check $\len t-J(t)<N_A$. A pair that passes all tests is a
mortal hole. If a mortal hole exists, \cref{thm:mo-short} supplies one among the
tested pairs.
\end{proof}
